\section{Master rate formulas at finite temperature}\label{app:D.25}

The rates of \cref{sec:8.2} follow from the golden rule with the interaction of \cref{eq:8.1-Hint} and a target in thermal equilibrium. Natural units are used; $\vect{q}$ is the momentum the dark matter loses and $\omega$ the energy the target gains.

\subsection{Scattering}

The dark matter is a plane wave of momentum $\vect{p} = \mchi\vect{v}$ in the box of volume $\Vol$ (\cref{app:A.3}), in the spin state $|s\rangle$, one of $2|\vect{S}_\chi| + 1$. The target is in the eigenstate $|i\rangle$ with the Gibbs weight $p_i$. By \cref{eq:8.1-Hint}, the matrix element to the dark-matter state $\vect{p}' = \vect{p} - \vect{q}$, $|s'\rangle$ and the target state $|f\rangle$ is
\begin{equation*}
  \mathcal{A}_{fi} = \frac1\Vol\sum_b\langle s'|\tilde V_b(\vect{q},\bar{\vect{v}}_\chi)|s\rangle\,\langle f|O_b(\vect{q})|i\rangle .
\end{equation*}
Energy conservation between the initial energy $E_i + E(\vect{p})$ and the final energy $E_f + E(\vect{p}')$ is $\delta(\omega_{\vect{q}} - E_f + E_i)$, with $\omega_{\vect{q}} = E(\vect{p}) - E(\vect{p}')$ the energy of \cref{eq:1.3-omegaq}, mass splitting included. The golden rule \cref{eq:2.1-golden}, averaged over the initial states as in \cref{eq:2.5-thermalrate}, and averaged over $s$ and summed over $s'$, gives
\begin{equation*}
  \varGamma = \sum_{\vect{p}'}\frac{2\pi}{2|\vect{S}_\chi| + 1}\sum_{s,s'}\sum_{i,f}p_i\,|\mathcal{A}_{fi}|^2\,\delta(\omega_{\vect{q}} - E_f + E_i) .
\end{equation*}
Squaring the sum over $b$ produces the pairs $(a, b)$, with the conjugate of the term $a$:
\begin{equation*}
  |\mathcal{A}_{fi}|^2 = \frac{1}{\Vol^2}\sum_{a,b}\langle s'|\tilde V_a|s\rangle^*\langle s'|\tilde V_b|s\rangle\;\langle i|O_a^\dagger|f\rangle\langle f|O_b|i\rangle .
\end{equation*}
Since $\langle s'|\tilde V_a|s\rangle^* = \langle s|\tilde V_a^\dagger|s'\rangle$, the sum over $s'$ is a completeness relation and the sum over $s$ a trace:
\begin{equation*}
  \frac{1}{2|\vect{S}_\chi| + 1}\sum_{s,s'}\langle s|\tilde V_a^\dagger|s'\rangle\langle s'|\tilde V_b|s\rangle = \frac{\Tr_\chi\bigl[\tilde V_a^\dagger\tilde V_b\bigr]}{2|\vect{S}_\chi| + 1} = \mathfrak{D}_{ab}(\vect{q},\vect{v}) ,
\end{equation*}
the tensor \cref{eq:8.1-D}. It depends on $\vect{v}$ and $\vect{q}$ through $\bar{\vect{v}}_\chi = \vect{v} - \vect{q}/2\mchi$. The final momenta are $\vect{p}' = \vect{p} - \vect{q}$, and the sum over them is $\sum_{\vect{p}'} = \Vol\int\dd^3q/(2\pi)^3$ (\cref{eq:2.1-box}). Collecting the factors,
\begin{equation*}
  \varGamma(\vect{v}) = \int\frac{\dd^3q}{(2\pi)^3}\sum_{a,b}\mathfrak{D}_{ab}(\vect{q},\vect{v})\;\frac{2\pi}{\Vol}\sum_{i,f}p_i\,\langle i|O_a^\dagger|f\rangle\langle f|O_b|i\rangle\,\delta(\omega_{\vect{q}} - E_f + E_i) ,
\end{equation*}
and the last factor is $S_{ab}(\vect{q},\omega_{\vect{q}})$ of \cref{eq:4.2-Sab}. This is \cref{eq:8.2-rate}. One power of $1/\Vol$ from the squared amplitude cancels against the sum over final momenta, the other is absorbed in $S_{ab}$, and the rate of one particle in the box is independent of $\Vol$.

The box holds $\rho_\chi \Vol/\mchi$ dark-matter particles with velocities distributed as $f_\chi(\vect{v})$, and a target mass $\rho_T\Vol$. The rate per unit target mass is then $(\rho_\chi/\mchi\rho_T)\int\dd^3v\,f_\chi(\vect{v})\,\varGamma(\vect{v})$. Inserting $1 = \int\dd\omega\,\delta(\omega - \omega_{\vect{q}})$ under the $\vect{q}$ integral and replacing $\omega_{\vect{q}}$ by $\omega$ in $S_{ab}$ gives the spectrum \cref{eq:8.2-spectrum}.

\subsection{Index order of the contraction}

The squared amplitude pairs $\tilde V_a^\dagger$ with $O_a^\dagger$, so the contraction is $\sum_{ab}\mathfrak{D}_{ab}S_{ab} = \Tr[\mathfrak{D}S^{\mathrm{T}}]$. Both matrices are Hermitian. Write each as a real symmetric part plus $\ii$ times a real antisymmetric part, $\mathfrak{D} = \mathfrak{D}_{\mathrm{R}} + \ii\mathfrak{D}_{\mathrm{I}}$ and $S = S_{\mathrm{R}} + \ii S_{\mathrm{I}}$. Then $S^{\mathrm{T}} = S_{\mathrm{R}} - \ii S_{\mathrm{I}}$, and since the trace of a symmetric matrix times an antisymmetric one vanishes,
\begin{equation*}
  \Tr[\mathfrak{D}S^{\mathrm{T}}] = \Tr[\mathfrak{D}_{\mathrm{R}}S_{\mathrm{R}}] + \Tr[\mathfrak{D}_{\mathrm{I}}S_{\mathrm{I}}], \qquad
  \Tr[\mathfrak{D}S] = \Tr[\mathfrak{D}_{\mathrm{R}}S_{\mathrm{R}}] - \Tr[\mathfrak{D}_{\mathrm{I}}S_{\mathrm{I}}] .
\end{equation*}
The two orders differ by $2\Tr[\mathfrak{D}_{\mathrm{I}}S_{\mathrm{I}}]$, which vanishes when either matrix is real. Both forms are real. Only $\Tr[\mathfrak{D}S^{\mathrm{T}}]$ is the rate, and only it is guaranteed nonnegative: $S^{\mathrm{T}} = S^*$ is positive semidefinite when $S$ is, and the trace of a product of two positive semidefinite matrices is nonnegative.

\subsection{Absorption from a quantized field}

Quantize a real scalar field in the box,
\begin{equation*}
  \phi(\vect{x}) = \sum_{\vect{k}}\frac{1}{\sqrt{2\omega_{\vect{k}}\Vol}}\bigl(a_{\vect{k}}\,\ee^{\ii\vect{k}\cdot\vect{x}} + a_{\vect{k}}^\dagger\,\ee^{-\ii\vect{k}\cdot\vect{x}}\bigr), \qquad \omega_{\vect{k}} = (m_\phi^2 + k^2)^{1/2} ,
\end{equation*}
with $[a_{\vect{k}}, a_{\vect{k}'}^\dagger] = \delta_{\vect{k}\vect{k}'}$; a vector field carries in addition a polarization vector for each mode. The coupling integrates the field against target bilinears, and the plane wave $\ee^{\ii\vect{k}\cdot\vect{x}}$ turns each bilinear into the row at $\vect{q} = \vect{k}$ (\cref{sec:8.1-absorption}); for the density, $\int\dd^3x\,\ee^{\ii\vect{k}\cdot\vect{x}}\op{\nd}_e(\vect{x}) = n_e(\vect{k})$. With $g^\phi_a$ the coefficient of row $a$ for the mode $\vect{k}$, derivatives of the field included,
\begin{equation*}
  H_{\mathrm{int}} = \sum_{\vect{k}}\frac{1}{\sqrt{2\omega_{\vect{k}}\Vol}}\Bigl[a_{\vect{k}}\sum_ag^\phi_a\,O_a(\vect{k}) + a_{\vect{k}}^\dagger\sum_ag^{\phi\,*}_a\,O_a(\vect{k})^\dagger\Bigr] .
\end{equation*}
From the state with $N_\phi$ quanta in mode $\vect{k}$ and the target in $|i\rangle$, absorption leads to $N_\phi - 1$ quanta and $|f\rangle$, with the matrix element $\sqrt{N_\phi/2\omega_{\vect{k}}\Vol}\,\sum_bg^\phi_b\langle f|O_b(\vect{k})|i\rangle$ and energy conservation $\delta(\omega_{\vect{k}} - E_f + E_i)$. The golden rule, thermally averaged and summed over $f$, gives
\begin{equation*}
  \varGamma_{N_\phi\to N_\phi-1} = \frac{N_\phi}{2\omega_{\vect{k}}}\sum_{a,b}g^{\phi\,*}_ag^\phi_b\;\frac{2\pi}{\Vol}\sum_{i,f}p_i\,\langle i|O_a^\dagger|f\rangle\langle f|O_b|i\rangle\,\delta(\omega_{\vect{k}} - E_f + E_i) = N_\phi\,\varGamma_{\mathrm{abs}} ,
\end{equation*}
with $\varGamma_{\mathrm{abs}}$ of \cref{eq:8.2-abs} at $\vect{p}_\phi = \vect{k}$. The rate is $N_\phi$ times the rate per quantum.

Emission leads to $N_\phi + 1$ quanta, with the matrix element $\sqrt{(N_\phi+1)/2\omega_{\vect{k}}\Vol}\,\sum_bg^{\phi\,*}_b\langle f|O_b^\dagger|i\rangle$ and $\delta(\omega_{\vect{k}} + E_f - E_i)$. Exchanging the names of $i$ and $f$, as in the proof of \cref{eq:2.5-db}, the weight becomes $p_f = p_i\,\ee^{-\beta\omega_{\vect{k}}}$ and the squared element becomes $|\sum_bg^\phi_b\langle f|O_b|i\rangle|^2$, so that $\varGamma_{N_\phi\to N_\phi+1} = (N_\phi + 1)\,\ee^{-\beta\omega_{\vect{k}}}\varGamma_{\mathrm{abs}}$. The field loses quanta at the net rate
\begin{equation*}
  \varGamma_{N_\phi\to N_\phi-1} - \varGamma_{N_\phi\to N_\phi+1} = N_\phi\,\varGamma_{\mathrm{abs}}\Bigl[1 - \Bigl(1 + \frac{1}{N_\phi}\Bigr)\ee^{-\beta\omega_{\vect{k}}}\Bigr] .
\end{equation*}
The term $1/N_\phi$ is spontaneous emission. For a halo field it is negligible, and per quantum the net rate is $\varGamma_{\mathrm{net}} = (1 - \ee^{-\beta\omega_{\vect{k}}})\varGamma_{\mathrm{abs}}$. By the fluctuation--dissipation theorem \cref{eq:4.2-fdt}, with $x_a = g^\phi_a$,
\begin{equation*}
  \begin{aligned}
  (1 - \ee^{-\beta\omega})\sum_{a,b}x_a^*x_bS_{ab} &= \ii\sum_{a,b}x_a^*x_b\bigl(R_{ab} - R_{ba}^*\bigr) \\
  &= \ii\Bigl[\sum_{a,b}x_a^*x_bR_{ab} - \Bigl(\sum_{a,b}x_a^*x_bR_{ab}\Bigr)^*\Bigr] = -2\,\Imag\sum_{a,b}x_a^*x_bR_{ab} ,
  \end{aligned}
\end{equation*}
where the middle step relabels $a \leftrightarrow b$ in the term with $R_{ba}^*$. Dividing by $2\omega_{\vect{k}}$ gives the second form of \cref{eq:8.2-abs}. The modes are occupied with $N_\phi$ quanta, and their number density $\sum_{\vect{k}}N_\phi/\Vol$ is $\rho_\chi/m_\phi$ up to corrections of order $v^2$, distributed in $\vect{k} = m_\phi\vect{v}$ as $f_\chi$. The event rate per unit target mass is $(\rho_\chi/m_\phi\rho_T)\langle\varGamma_{\mathrm{abs}}\rangle$.

\subsection{Absorption from a classical field}

For $N_\phi \gg 1$ the field can be treated as classical. Write it as a sum of modes with complex amplitudes,
\begin{equation*}
  \phi(\vect{x},t) = \sum_{\vect{k}}\bigl[\phi_{\vect{k}}\,\ee^{\ii(\vect{k}\cdot\vect{x} - \omega_{\vect{k}}t)} + \phi_{\vect{k}}^*\,\ee^{-\ii(\vect{k}\cdot\vect{x} - \omega_{\vect{k}}t)}\bigr] .
\end{equation*}
A virialized halo is a superposition of many such waves with independent, uniformly distributed phases. By the central limit theorem the field at any point is then a Gaussian random variable, and all its correlations follow from $\langle\phi_{\vect{k}}\phi_{\vect{k}'}^*\rangle = \delta_{\vect{k}\vect{k}'}\langle|\phi_{\vect{k}}|^2\rangle$, with $\langle\phi_{\vect{k}}\phi_{\vect{k}'}\rangle = 0$. The energy density, $\frac12\dot\phi^2 + \frac12(\nabla\phi)^2 + \frac12m_\phi^2\phi^2$, has the average $\sum_{\vect{k}}(\omega_{\vect{k}}^2 + k^2 + m_\phi^2)\langle|\phi_{\vect{k}}|^2\rangle = \sum_{\vect{k}}2\omega_{\vect{k}}^2\langle|\phi_{\vect{k}}|^2\rangle$. The number density of quanta in mode $\vect{k}$ is its energy density divided by $\omega_{\vect{k}}$,
\begin{equation*}
  \nd_{\vect{k}} = 2\omega_{\vect{k}}\,\langle|\phi_{\vect{k}}|^2\rangle , \qquad \sum_{\vect{k}}\nd_{\vect{k}} \simeq \frac{\rho_\chi}{m_\phi} .
\end{equation*}
For a single mode, $\phi = \phi_0\cos(m_\phi t)$ has $|\phi_{\vect{0}}| = \phi_0/2$ and $\rho_\chi = \frac12m_\phi^2\phi_0^2$, as in \cref{sec:2.1}.

The interaction is now a sum of harmonic perturbations,
\begin{equation*}
  H'(t) = \sum_{\vect{k}}\bigl[\ee^{-\ii\omega_{\vect{k}}t}\,\phi_{\vect{k}}\,\op{O}_{\vect{k}} + \ee^{\ii\omega_{\vect{k}}t}\,\phi_{\vect{k}}^*\,\op{O}_{\vect{k}}^\dagger\bigr], \qquad
  \op{O}_{\vect{k}} = \sum_ag^\phi_a\,O_a(\vect{k}) .
\end{equation*}
The first-order amplitude of \cref{sec:2.1} is a sum over modes, and its square contains products $\phi_{\vect{k}}\phi_{\vect{k}'}^*$. The Gaussian field average sets those with $\vect{k} \ne \vect{k}'$ to zero; a single realization does the same in time, once the observation time exceeds the inverse spread of the frequencies $\omega_{\vect{k}}$, which is the coherence time $t_{\mathrm{coh}}$. What remains is the sum over modes of the harmonic rate \cref{eq:2.1-harmonic}, with $\op{O} = \phi_{\vect{k}}\op{O}_{\vect{k}}$. Summed over $f$ and thermally averaged, the absorption term gives the rate in the box
\begin{equation*}
  \begin{aligned}
  \sum_{\vect{k}}2\pi\,\langle|\phi_{\vect{k}}|^2\rangle\sum_{i,f}p_i\,|\langle f|\op{O}_{\vect{k}}|i\rangle|^2\,\delta(\omega_{\vect{k}} - E_f + E_i)
  &= \Vol\sum_{\vect{k}}\frac{\nd_{\vect{k}}}{2\omega_{\vect{k}}}\sum_{a,b}g^{\phi\,*}_ag^\phi_b\,S_{ab}(\vect{k},\omega_{\vect{k}}) \\
  &= \Vol\sum_{\vect{k}}\nd_{\vect{k}}\,\varGamma_{\mathrm{abs}}(\vect{k}) .
  \end{aligned}
\end{equation*}
Each quantum per unit volume is absorbed at the rate $\varGamma_{\mathrm{abs}}$, as in the quantized calculation. The emission term gives $\ee^{-\beta\omega_{\vect{k}}}$ times the same, with no spontaneous part. The two descriptions agree for $N_\phi \gg 1$. The classical one shows what the golden rule needs: $t \gg t_{\mathrm{coh}}$, so that the modes add incoherently, and $t$ long compared with the inverse width of the target's response, so that \cref{eq:2.1-harmonic} holds for each mode. With $t_{\mathrm{coh}} \sim 1/(m_\phi v_0^2)$ this is about \SI{1}{ns} at $m_\phi = \SI{1}{eV}$ and grows as $1/m_\phi$.

\subsection{In-medium mixing of the dark photon}

A dark photon $A'$ couples to the charges of the target with $\zeta e$ in place of $e$, so the charges respond to the vector potential $\zeta\vect{A}' + \vect{A}$, where $\vect{A}$ is the ordinary electromagnetic field that the target itself produces. Take a transverse mode of frequency $\omega$ and wavevector $\vect{k}$ with $\omega^2 - k^2 = m_\phi^2$. The induced current is fixed by the photon self-energy of \cref{sec:4.4}, $\vect{J} = -\Pi_{\mathrm{em}}(\zeta\vect{A}' + \vect{A})$, the relation whose homogeneous case is the medium's wave equation $(\omega^2 - k^2 - \Pi_{\mathrm{em}})\vect{A} = 0$. The current sources the ordinary field through $(k^2 - \omega^2)\vect{A} = \vect{J}$, so that $\vect{A} = -\vect{J}/m_\phi^2$ at this frequency and wavevector. Eliminating $\vect{A}$,
\begin{equation*}
  \vect{J} = -\zeta\,\frac{m_\phi^2\,\Pi_{\mathrm{em}}}{m_\phi^2 - \Pi_{\mathrm{em}}}\,\vect{A}' .
\end{equation*}
The dark photon therefore sees the self-energy $\zeta^2m_\phi^2\Pi_{\mathrm{em}}/(m_\phi^2 - \Pi_{\mathrm{em}})$ instead of $\zeta^2\Pi_{\mathrm{em}}$, and its absorption rate is $-1/\omega$ times the imaginary part. Since $\Imag[\Pi_{\mathrm{em}}(m_\phi^2 - \Pi_{\mathrm{em}}^*)] = m_\phi^2\Imag\Pi_{\mathrm{em}}$,
\begin{equation*}
  \varGamma_{\mathrm{net}} = -\frac{\zeta^2}{\omega}\,\frac{m_\phi^4\,\Imag\Pi_{\mathrm{em}}}{|m_\phi^2 - \Pi_{\mathrm{em}}|^2} = \zeta_{\mathrm{eff}}^2\,\sigma_1(\omega) ,
\end{equation*}
with $\Imag\Pi_{\mathrm{em}} = -\omega\sigma_1$ from $\Pi_{\mathrm{em}} = -\ii\omega\sigma_{\mathrm{cond}}$ and $\zeta_{\mathrm{eff}}$ of \cref{sec:1.4}~\cites[Eqs.~(5) and (6)]{hochberg2017absorption}[Eqs.~(20)--(22)]{mitridate2021dark}. To first order in $\Pi_{\mathrm{em}}/m_\phi^2$, the factor is one and the rate is $\zeta^2\sigma_1$, the result of the proper response alone. The photon-reducible part matters when $m_\phi$ is comparable to the plasma frequency, \SI{16.6}{eV} in silicon (\cref{sec:4.3}).

\subsection{Detailed balance for the dark matter}

Let the dark matter go from $\vect{p}$ to one final momentum $\vect{p}'$, depositing $\omega = E(\vect{p}) - E(\vect{p}')$, where the two states may belong to different mass eigenstates $\chi_1$ and $\chi_2$ of equal spin. By the golden rule and the steps that led to \cref{eq:8.2-rate}, the rate into $\vect{p}'$ is
\begin{equation*}
  \varGamma(\vect{p}\to\vect{p}') = \frac{2\pi}{\Vol^2}\sum_{a,b}\mathfrak{D}_{ab}\sum_{i,f}p_i\,\langle i|O_a^\dagger|f\rangle\langle f|O_b|i\rangle\,\delta(\omega - E_f + E_i) .
\end{equation*}
The reverse transition, from $\vect{p}'$ and the target state $|f\rangle$ to $\vect{p}$ and $|i\rangle$, has the matrix element of the Hermitian interaction between the exchanged states, the complex conjugate of the forward one:
\begin{equation*}
  \langle\vect{p}, s; i|H_{\mathrm{int}}|\vect{p}', s'; f\rangle = \frac1\Vol\sum_b\langle s|\tilde V_b^\dagger|s'\rangle\,\langle i|O_b^\dagger|f\rangle .
\end{equation*}
Squaring, averaging over $s'$ and summing over $s$ gives
\begin{equation*}
  \sum_{a,b}\frac{\Tr_\chi[\tilde V_a\tilde V_b^\dagger]}{2|\vect{S}_\chi| + 1}\,\langle f|O_a|i\rangle\langle i|O_b^\dagger|f\rangle ,
\end{equation*}
and the cyclic trace is $\Tr_\chi[\tilde V_b^\dagger\tilde V_a] = (2|\vect{S}_\chi| + 1)\,\mathfrak{D}_{ba}$. The initial target state is now $|f\rangle$, with weight $p_f$, and energy conservation is $\delta(-\omega - E_i + E_f)$, the same delta function. Relabeling $a \leftrightarrow b$,
\begin{equation*}
  \varGamma(\vect{p}'\to\vect{p}) = \frac{2\pi}{\Vol^2}\sum_{a,b}\mathfrak{D}_{ab}\sum_{i,f}p_f\,\langle i|O_a^\dagger|f\rangle\langle f|O_b|i\rangle\,\delta(\omega - E_f + E_i) ,
\end{equation*}
and on the support of the delta function $p_f = p_i\,\ee^{-\beta\omega}$. This is \cref{eq:8.2-db}. Neither the target's symmetries nor the form of the coupling entered; only the Hermiticity of the interaction and the Gibbs weights did. If the dark-matter states had different spins, the ratio would carry in addition the ratio of their multiplicities from the spin averages.
